\uuid{7Z2k}
\titre{ Sous-espaces vectoriels de $\mathbb{R}^4$ }

\niveau{L1} 				%L1, L2, L3, MPSI, MP, PCSI, PC, PSI...
\module{ Espaces vectoriels } 	%Analyse, Algèbre...
\chapitre{Bases et espaces supplémentaires}   			%Continuité, Groupes, Fonctions de plusieurs variables...
\sousChapitre{}			%Optimisation, Diagonalisation d'une matrice, Calcul de dérivées partielles...

\theme{}				%Fonction de répartition, Division euclidienne de polynômes, ...
\auteur{Quentin Liard}
\datecreate{ 2026-02-19}
\organisation{AMSCC}			%AMSCC, Exo7, ...
\difficulte{2}			%1, 2, 3, 4 ou 5

\contenu{

\texte{ 
On se place dans $\mathbb{R}^4$.
On considère
$E=\{(x,y,z,t)\in\mathbb{R}^4\mid x+y+z+t=0\}$
et \\
$F=\{(x,y,z,t)\in\mathbb{R}^4\mid x-y=0\ \text{et}\ z-t=0\}$.
}

\begin{enumerate}
\item   \question{Montrer que $E$ et $F$ sont des sous-espaces vectoriels de $\mathbb{R}^4$.}
\indication{}
\reponse{\textbf{Pour $E$ :}  
On a $E=\{(x,y,z,t)\mid x+y+z+t=0\}$.

\begin{itemize}
\item Le vecteur nul $(0,0,0,0)$ vérifie $0+0+0+0=0$, donc $0\in E$.
\item Si $u=(x,y,z,t)\in E$ et $v=(x',y',z',t')\in E$, alors
\[
(x+x')+(y+y')+(z+z')+(t+t')=(x+y+z+t)+(x'+y'+z'+t')=0.
\]
Donc $u+v\in E$.
\item Pour $\lambda\in\mathbb{R}$,
\[
\lambda(x+y+z+t)=\lambda\cdot 0=0,
\]
donc $\lambda u\in E$.
\end{itemize}

Ainsi $E$ est un sous-espace vectoriel.

\medskip

\textbf{Pour $F$ :}  
On a $F=\{(x,y,z,t)\mid x-y=0 \text{ et } z-t=0\}$.

\begin{itemize}
\item $(0,0,0,0)$ vérifie les deux équations, donc $0\in F$.
\item Si $u,v\in F$, alors
\[
(x+x')-(y+y')=(x-y)+(x'-y')=0
\]
et
\[
(z+z')-(t+t')=(z-t)+(z'-t')=0,
\]
donc $u+v\in F$.
\item Pour $\lambda\in\mathbb{R}$,
\[
\lambda(x-y)=0 \quad \text{et} \quad \lambda(z-t)=0,
\]
donc $\lambda u\in F$.
\end{itemize}

Ainsi $F$ est un sous-espace vectoriel de $\mathbb{R}^4$.}

\item   \question{Donner une base de $E$ et déterminer $\dim E$.}
\indication{}
\reponse{Paramétrisation : $x=-y-z-t$ donc $(x,y,z,t)=y(-1,1,0,0)+z(-1,0,1,0)+t(-1,0,0,1)$.
Une base est $\{(-1,1,0,0),(-1,0,1,0),(-1,0,0,1)\}$ car les vecteurs sont libres . En effet soit $\lambda_1,\lambda_2,\lambda_3 \in \mathbb{R}$ tels que $$\lambda_1 (-1,1,0,0)+\lambda_2 (-1,0,1,0)+\lambda_3 (-1,0,0,1)=0.$$ Ceci implique $$-\lambda_1-\lambda_2-\lambda_3=0,\quad \lambda_1=0,\quad \lambda_2=0,\quad \lambda_3=0,$$
Ainsi $\dim E=3$.}

\item   \question{Donner une base de $F$ et déterminer $\dim F$.}
\indication{}
\reponse{Les contraintes sont $x=y$ et $z=t$. Posons . Alors $(x,y,z,t)=(x,x,b,b)=x(1,1,0,0)+z(0,0,1,1)$.
Donc $\dim F=2$ et une base est $\{(1,1,0,0),(0,0,1,1)\}$.}

\item   \question{Déterminer $E\cap F$.}
\indication{}
\reponse{Dans $F$, un vecteur est $(a,a,b,b)$. Il est dans $E$ ssi $a+a+b+b=0$, soit $2a+2b=0$, donc $a=-b$.
Ainsi $E\cap F=\{(a,a,-a,-a)\mid a\in\mathbb{R}\}=\mathrm{Vect}((1,1,-1,-1))$ et $\dim(E\cap F)=1$.}

\item   \question{A-t-on $\mathbb{R}^4=E+F$ ?}
\indication{}
\reponse{Oui, car $\dim(E+F)=\dim E+\dim F-\dim(E\cap F)=3+2-1=4$, donc $E+F$ est un sous-espace de dimension 4, donc $E+F=\mathbb{R}^4$.}

\item   \question{A-t-on $\mathbb{R}^4=E\oplus F$ ?}
\indication{}
\reponse{Non, une somme directe impose $E\cap F=\{0\}$. Or $E\cap F=\mathrm{Vect}((1,1,-1,-1))\neq\{0\}$.}

\end{enumerate}

}